\documentclass[12pt]{article}

% ==================== بسته‌ها ====================
\usepackage{xepersian}
\settextfont{IRANSansX} % فونت فارسی نصب‌شده روی این سیستم؛ در صورت نیاز به فونت دیگر تغییر دهید

\usepackage{amsmath}
\usepackage{amssymb}
\usepackage{graphicx}
\usepackage{booktabs}
\usepackage{array}
\usepackage{longtable}
\usepackage{caption}
\usepackage{geometry}
\usepackage{tikz}
\usepackage{xcolor}

\geometry{margin=2.7cm}

% ==================== جای خالی برای اسکرین‌شات ====================
% وقتی اسکرین‌شات مدار را گرفتید، این دستور را با دستور زیر جایگزین کنید:
%   \includegraphics[width=0.9\textwidth]{images/NAME.png}
\newcommand{\screenshotplaceholder}[1]{%
  \begin{center}
    \setlength{\fboxsep}{0pt}
    \fbox{%
      \begin{minipage}[c][6.5cm][c]{0.85\textwidth}
        \centering
        {\Large جای خالی برای اسکرین‌شات مدار}\\[4mm]
        {\normalsize #1}
      \end{minipage}%
    }
  \end{center}
}

\title{
  {\Large دانشکده مهندسی کامپیوتر، دانشگاه صنعتی شریف}\\[2mm]
  {\large درس آزمایشگاه معماری کامپیوتر (CE103)}\\[8mm]
  \rule{\linewidth}{0.6pt}\\[4mm]
  {\huge \bfseries گزارش‌کار آزمایش اول}\\[2mm]
  {\Large طراحی و پیاده‌سازی جمع‌کننده دهدهی سه‌رقمی (BCD Adder)}\\[2mm]
  \rule{\linewidth}{0.6pt}
}
\author{نام و نام‌خانوادگی: \dotfill \\ شماره دانشجویی: \dotfill}
\date{استاد درس: دکتر سربازی \hspace{1cm} تاریخ: \dotfill}

\begin{document}

\maketitle
\thispagestyle{empty}

\begin{abstract}
در این آزمایش یک جمع‌کننده‌ی دهدهی سه‌رقمی (\lr{3-digit BCD Adder}) در نرم‌افزار \lr{Logisim Evolution} طراحی و پیاده‌سازی شده است. مدار به‌صورت سلسله‌مراتبی و از پایین به بالا ساخته شده: یک تمام‌جمع‌کننده‌ی ۱ بیتی (\lr{full\_adder})، که پایه‌ی ساخت یک جمع‌کننده‌ی یک‌رقمی دهدهی با تصحیح خودکار \lr{BCD} (\lr{bcd\_digit}) است، و در نهایت سه نمونه از این جمع‌کننده‌ی رقمی برای ساخت یک جمع‌کننده‌ی سه‌رقمی کامل (\lr{main}) زنجیر شده‌اند. مدار به‌طور کامل با بردارهای تست تأیید شده است.
\end{abstract}

\tableofcontents
\newpage

% ==================================================================
\section{هدف آزمایش}
% ==================================================================

هدف از این آزمایش، آشنایی عملی با طراحی مدارهای منطقی ترکیبی سلسله‌مراتبی و پیاده‌سازی یک جمع‌کننده‌ی اعداد دهدهی کدشده به‌صورت \lr{BCD} (\lr{Binary Coded Decimal}) است. برخلاف جمع باینری معمولی، جمع دو رقم \lr{BCD} به یک مرحله‌ی تصحیح اضافه نیاز دارد تا نتیجه در بازه‌ی مجاز \lr{(۰ تا ۹)} باقی بماند. اهداف مشخص این آزمایش عبارت‌اند از:

\begin{itemize}
  \item پیاده‌سازی یک تمام‌جمع‌کننده‌ی ۱ بیتی (\lr{Full Adder}) با گیت‌های پایه.
  \item طراحی مدار تصحیح \lr{BCD} برای جمع دو رقم دهدهی به‌همراه رقم نقلی ورودی.
  \item ساخت یک جمع‌کننده‌ی سه‌رقمی با زنجیر کردن سه نمونه از جمع‌کننده‌ی یک‌رقمی.
  \item صحت‌سنجی مدار با بردارهای تست جامع (\lr{exhaustive test vectors}).
\end{itemize}

% ==================================================================
\section{مبانی تئوری}
% ==================================================================

\subsection{کدگذاری \lr{BCD}}

در کدگذاری \lr{BCD}، هر رقم دهدهی (۰ تا ۹) با ۴ بیت باینری نمایش داده می‌شود. برخلاف نمایش باینری خالص، در \lr{BCD} کدهای ۱۰۱۰ تا ۱۱۱۱ (معادل ۱۰ تا ۱۵) بلااستفاده و نامعتبر هستند  -  هر رقم باید همیشه بین $0000_2$ تا $1001_2$ باشد.

\subsection{مسئله‌ی جمع دو رقم \lr{BCD}}

اگر دو رقم \lr{BCD} (هرکدام حداکثر ۹) به‌همراه یک رقم نقلی ورودی (حداکثر ۱) با یک جمع‌کننده‌ی باینری معمولی ۴ بیتی جمع زده شوند، نتیجه‌ی خام می‌تواند تا:
\[
9 + 9 + 1 = 19
\]
برسد که در بازه‌ی معتبر \lr{BCD} (۰ تا ۹) نیست. اگر جمع خام از ۹ بیشتر باشد، باید مقدار ۶ ($0110_2$) به آن اضافه شود تا:
\begin{enumerate}
  \item رقم باقی‌مانده در بازه‌ی صحیح \lr{BCD} قرار بگیرد، و
  \item یک رقم نقلی دهدهی (نه باینری) به رقم بعدی منتقل شود.
\end{enumerate}

\subsection{چرا رقم نقلی باینری کافی نیست؟}

رقم نقلی خروجی یک جمع‌کننده‌ی ۴ بیتی معمولی، تنها زمانی ۱ می‌شود که جمع خام به ۱۶ یا بیشتر برسد (سرریز باینری). اما مقادیر بین ۱۰ تا ۱۵ نیز خارج از بازه‌ی \lr{BCD} هستند، در حالی که هیچ سرریز باینری‌ای رخ نداده است. بنابراین باید یک مدار تشخیص جداگانه طراحی شود که وضعیت «$\text{جمع خام} > 9$» را به‌درستی تشخیص دهد.

% ==================================================================
\section{طراحی و پیاده‌سازی مدار}
% ==================================================================

مدار در سه سطح، از پایین به بالا، طراحی شده است:

\begin{center}
\begin{tabular}{@{}lll@{}}
\toprule
\textbf{مدار} & \textbf{نقش} & \textbf{ساخته‌شده از} \\
\midrule
\lr{full\_adder} & تمام‌جمع‌کننده‌ی ۱ بیتی & گیت‌های \lr{XOR}, \lr{AND}, \lr{OR} \\
\lr{bcd\_digit} & جمع‌کننده‌ی یک‌رقمی دهدهی با تصحیح & ۸ نمونه \lr{full\_adder} \\
\lr{main} & جمع‌کننده‌ی سه‌رقمی کامل & ۳ نمونه \lr{bcd\_digit} \\
\bottomrule
\end{tabular}
\end{center}

نکته‌ی مهم طراحی: در نسخه‌ی نهایی، هیچ ماژول واسط اضافه‌ای (مثل یک جمع‌کننده‌ی ۴بیتی جدا) وجود ندارد؛ \lr{bcd\_digit} مستقیما از ۸ نمونه‌ی \lr{full\_adder} ساخته شده تا تعداد ماژول‌ها به حداقل (۳ ماژول) برسد.

% ------------------------------------------------------------------
\subsection{تمام‌جمع‌کننده‌ی ۱ بیتی  -  \lr{full\_adder}}
% ------------------------------------------------------------------

\subsubsection{پایه‌ها}

\begin{center}
\begin{tabular}{@{}lll@{}}
\toprule
\textbf{پایه} & \textbf{جهت} & \textbf{پهنا} \\
\midrule
\lr{A}, \lr{B}, \lr{Cin} & ورودی & ۱ بیت \\
\lr{Sum}, \lr{Cout} & خروجی & ۱ بیت \\
\bottomrule
\end{tabular}
\end{center}

\subsubsection{منطق پیاده‌سازی}

مدار با ۲ گیت \lr{XOR}، ۲ گیت \lr{AND} و ۱ گیت \lr{OR} ساخته شده است:

\begin{align}
X_1 &= A \oplus B \qquad &&\text{(نیمه‌جمع، بدون در نظر گرفتن رقم نقلی ورودی)}\\
P_1 &= A \land B \qquad &&\text{(آیا از جمع } A,B \text{ به‌تنهایی رقم نقلی تولید می‌شود؟)}\\
\text{Sum} &= X_1 \oplus C_{in} \qquad &&\text{(جمع نهایی سه ورودی)}\\
P_2 &= X_1 \land C_{in} \qquad &&\text{(آیا از جمع } X_1, C_{in} \text{ رقم نقلی تولید می‌شود؟)}\\
\text{Cout} &= P_1 \lor P_2 \qquad &&\text{(رقم نقلی خروجی نهایی)}
\end{align}

به‌عبارت شهودی، \lr{Cout} زمانی ۱ می‌شود که \textbf{حداقل دو تا} از سه ورودی \lr{A}, \lr{B}, \lr{Cin} برابر ۱ باشند (یک رأی‌گیری اکثریت میان سه بیت)  -  که دقیقا همان چیزی است که $P_1 \lor P_2$ محاسبه می‌کند.

\subsubsection{جدول درستی}

\begin{center}
\begin{tabular}{@{}ccc|cc@{}}
\toprule
\lr{A} & \lr{B} & \lr{Cin} & \lr{Sum} & \lr{Cout} \\
\midrule
0&0&0&0&0\\
0&0&1&1&0\\
0&1&0&1&0\\
0&1&1&0&1\\
1&0&0&1&0\\
1&0&1&0&1\\
1&1&0&0&1\\
1&1&1&1&1\\
\bottomrule
\end{tabular}
\end{center}

\screenshotplaceholder{نمای داخلی مدار \lr{full\_adder} در \lr{Logisim} (۲ گیت \lr{XOR}، ۲ گیت \lr{AND}، ۱ گیت \lr{OR})}

% ------------------------------------------------------------------
\subsection{جمع‌کننده‌ی یک‌رقمی دهدهی با تصحیح  -  \lr{bcd\_digit}}
% ------------------------------------------------------------------

\subsubsection{پایه‌ها}

\begin{center}
\begin{tabular}{@{}lll@{}}
\toprule
\textbf{پایه} & \textbf{جهت} & \textbf{پهنا} \\
\midrule
\lr{A}, \lr{B} & ورودی & ۴ بیت (رقم \lr{BCD}، ۰ تا ۹) \\
\lr{Cin} & ورودی & ۱ بیت \\
\lr{S} & خروجی & ۴ بیت (رقم نتیجه، ۰ تا ۹) \\
\lr{Cout} & خروجی & ۱ بیت (رقم نقلی دهدهی) \\
\bottomrule
\end{tabular}
\end{center}

این مدار در سه مرحله کار می‌کند که در ادامه هرکدام به‌طور کامل توضیح داده می‌شوند.

\paragraph{مرحله‌ی اول: جمع باینری خام.}
ورودی‌های \lr{A} و \lr{B} با یک تفکیک‌کننده (\lr{Splitter}) به بیت‌های تکی $A_0..A_3$ و $B_0..B_3$ شکسته می‌شوند و در زنجیره‌ی اول از ۴ تمام‌جمع‌کننده (\lr{FA0} تا \lr{FA3}) به‌صورت \lr{ripple-carry} جمع زده می‌شوند:

\begin{center}
\begin{tabular}{@{}lllll@{}}
\toprule
\textbf{نمونه} & \textbf{بیت} & \textbf{ورودی‌ها} & \textbf{خروجی جمع} & \textbf{رقم نقلی} \\
\midrule
\lr{FA0} & ۰ (کم‌ارزش) & $A_0, B_0, C_{in}$ & $S_0$ & $CY_0$ \\
\lr{FA1} & ۱ & $A_1, B_1, CY_0$ & $S_1$ & $CY_1$ \\
\lr{FA2} & ۲ & $A_2, B_2, CY_1$ & $S_2$ & $CY_2$ \\
\lr{FA3} & ۳ (پرارزش) & $A_3, B_3, CY_2$ & $S_3$ & $C_{outH}$ \\
\bottomrule
\end{tabular}
\end{center}

نتیجه‌ی این مرحله عدد ۴ بیتی $S_3 S_2 S_1 S_0$ (جمع باینری خام، \textbf{هنوز تصحیح‌نشده}) به‌همراه $C_{outH}$ (رقم نقلی خروجی این جمع ۴ بیتی، معادل وزن ۱۶) است.

\paragraph{مرحله‌ی دوم: تشخیص نیاز به تصحیح.}
باید بررسی شود که آیا عدد $S_3S_2S_1S_0$ (با احتساب $C_{outH}$) از ۹ بزرگ‌تر است یا نه. چون $C_{outH}$ به‌تنهایی فقط سرریز $\geq 16$ را نشان می‌دهد، فرمول کمینه‌ی زیر (استاندارد در طراحی جمع‌کننده‌ی \lr{BCD}) برای پوشش کامل بازه‌ی ۱۰ تا ۱۹ استفاده شده است:

\begin{equation}
\text{CORR} = C_{outH} \lor (S_3 \land S_2) \lor (S_3 \land S_1)
\label{eq:corr}
\end{equation}

این فرمول با تنها ۳ گیت پیاده‌سازی شده است:

\begin{center}
\begin{tabular}{@{}llll@{}}
\toprule
\textbf{سیگنال} & \textbf{گیت} & \textbf{فرمول} & \textbf{معنی} \\
\midrule
\lr{G1OUT} & \lr{OR} & $S_2 \lor S_1$ & آیا یکی از بیت‌های ۱ یا ۲ روشن است؟ \\
\lr{G2OUT} & \lr{AND} & $S_3 \land \text{G1OUT}$ & آیا بیت ۳ \emph{و} یکی از بیت‌های ۱/۲ هم‌زمان روشن‌اند؟ \\
\lr{CORR} & \lr{OR} & $C_{outH} \lor \text{G2OUT}$ & سیگنال نهایی تصحیح \\
\bottomrule
\end{tabular}
\end{center}

سیگنال \lr{CORR} مستقیما به پایه‌ی خروجی \lr{Cout} این مدار وصل است؛ یعنی \textbf{رقم نقلی دهدهی، همان سیگنال تصحیح است.}

\paragraph{چرا این فرمول درست کار می‌کند؟}
اگر $C_{outH}=1$، جمع خام حتما $\geq 16 > 9$ است. اگر $C_{outH}=0$، جمع خام حداکثر ۱۵ است؛ در این حالت باید $S_3=1$ باشد تا مقدار $\geq 8$ شود. با $S_3=1$: اگر $S_2=1$ مقدار حداقل ۱۲ است (همیشه $>9$)؛ اگر $S_2=0$ ولی $S_1=1$ مقدار حداقل ۱۰ است (باز هم $>9$). این دقیقا همان چیزی است که عبارت $(S_3\land S_2)\lor(S_3\land S_1)$ تشخیص می‌دهد.

\paragraph{مرحله‌ی سوم: جمع تصحیحی $+6$.}
اگر \lr{CORR}$=1$ باشد، باید مقدار ۶ (باینری $0110$) به عدد خام اضافه شود. چون بیت‌های ۰ و ۳ عدد ۶ صفر هستند، کافی است \lr{CORR} فقط به بیت‌های ۱ و ۲ اضافه شود  -  بدون نیاز به جمع‌کننده یا گیت \lr{AND} ۴ بیتی جداگانه. این کار با زنجیره‌ی دوم از ۴ تمام‌جمع‌کننده (\lr{FAC0} تا \lr{FAC3}) انجام می‌شود:

\begin{center}
\begin{tabular}{@{}lllllll@{}}
\toprule
\textbf{نمونه} & \textbf{بیت} & \textbf{ورودی} \lr{A} & \textbf{ورودی} \lr{B} & \lr{Cin} & \textbf{جمع نهایی} & \textbf{رقم نقلی} \\
\midrule
\lr{FAC0} & ۰ & $S_0$ (خام) & \lr{ZERO} (۰) & \lr{ZERO} & $SD_0$ & $CYC_0$ \\
\lr{FAC1} & ۱ & $S_1$ (خام) & \lr{CORR} & $CYC_0$ & $SD_1$ & $CYC_1$ \\
\lr{FAC2} & ۲ & $S_2$ (خام) & \lr{CORR} & $CYC_1$ & $SD_2$ & $CYC_2$ \\
\lr{FAC3} & ۳ & $S_3$ (خام) & \lr{ZERO} & $CYC_2$ & $SD_3$ & (استفاده نمی‌شود) \\
\bottomrule
\end{tabular}
\end{center}

بیت‌های $SD_0..SD_3$ (رقم نهایی و تصحیح‌شده) با یک تفکیک‌کننده‌ی ترکیبی به پایه‌ی خروجی \lr{S} وصل می‌شوند و هم‌زمان برای نمایش بصری به یک نمایشگر هفت‌قطعه‌ای (\lr{Hex Digit Display}) می‌روند.

\emph{چرا رقم نقلی \lr{FAC3} استفاده نمی‌شود؟} چون حداکثر مقدار ممکن در این جمع، $9+6=15$ است که همیشه در ۴ بیت جا می‌شود؛ بنابراین این جمع هرگز سرریز نمی‌کند.

\screenshotplaceholder{نمای کلی مدار \lr{bcd\_digit}: زنجیره‌ی اول \lr{full\_adder} (جمع خام)، گیت‌های تشخیص تصحیح، و زنجیره‌ی دوم \lr{full\_adder} (جمع $+6$)}

\subsubsection{مثال عددی گام‌به‌گام}
جمع $7+8$ (با $C_{in}=0$) را در نظر بگیرید:
\begin{itemize}
  \item جمع خام: $0111 + 1000 = 1111$ یعنی $S_3S_2S_1S_0=1111$ و $C_{outH}=0$.
  \item تشخیص: $\text{G1OUT}=S_2\lor S_1=1\lor1=1$؛ $\text{G2OUT}=S_3\land\text{G1OUT}=1\land1=1$؛ $\text{CORR}=C_{outH}\lor\text{G2OUT}=0\lor1=1$.
  \item تصحیح: $1111 + 0110 = 1\,0101$؛ چون رقم نقلی این جمع در ۴ بیت جا نمی‌شود (نتیجه ۵ بیتی است ولی همان‌طور که گفته شد این رقم نقلی نادیده گرفته می‌شود چون از قبل توسط \lr{CORR} محاسبه شده)، ۴ بیت پایین یعنی $0101=5$ نتیجه‌ی رقم است.
  \item نتیجه: $S=5$, $\text{Cout}=1$  -  که دقیقا معادل $7+8=15$ در پایه‌ی ده است (رقم ۵ با یک رقم نقلی به مرتبه‌ی بعد).
\end{itemize}

% ------------------------------------------------------------------
\subsection{مدار اصلی  -  \lr{main} (سه رقم زنجیره‌شده)}
% ------------------------------------------------------------------

\subsubsection{پایه‌ها}

\begin{center}
\begin{tabular}{@{}lll@{}}
\toprule
\textbf{پایه} & \textbf{جهت} & \textbf{پهنا} \\
\midrule
\lr{A100}, \lr{A10}, \lr{A1} & ورودی & ۴ بیت هرکدام (رقم صدگان، دهگان، یکان عدد اول) \\
\lr{B100}, \lr{B10}, \lr{B1} & ورودی & ۴ بیت هرکدام (رقم صدگان، دهگان، یکان عدد دوم) \\
\lr{Cin} & ورودی & ۱ بیت \\
\lr{S100}, \lr{S10}, \lr{S1} & خروجی & ۴ بیت هرکدام \\
\lr{Cout} & خروجی & ۱ بیت \\
\bottomrule
\end{tabular}
\end{center}

\subsubsection{معماری}

سه نمونه از \lr{bcd\_digit}  -  به نام‌های \lr{digit\_1}، \lr{digit\_10} و \lr{digit\_100}  -  دقیقا مانند جمع دستی چند رقمی روی کاغذ، به‌صورت زنجیره‌ای به هم وصل شده‌اند:

\begin{center}
\begin{tabular}{@{}llllll@{}}
\toprule
\textbf{نمونه} & \textbf{رقم} & \lr{A} & \lr{B} & \lr{Cin} & \lr{Cout} می‌رود به \\
\midrule
\lr{digit\_1} & یکان & \lr{A1} & \lr{B1} & \lr{Cin} (ورودی مدار) & تونل \lr{Cout1} \\
\lr{digit\_10} & دهگان & \lr{A10} & \lr{B10} & تونل \lr{Cout1} & تونل \lr{Cout10} \\
\lr{digit\_100} & صدگان & \lr{A100} & \lr{B100} & تونل \lr{Cout10} & پایه‌ی \lr{Cout} مدار \\
\bottomrule
\end{tabular}
\end{center}

هر رقم روی یک نمایشگر \lr{Hex Digit Display} جداگانه نشان داده می‌شود و \lr{Cout} نهایی به یک \lr{LED} وصل است (در صورت سرریز از عدد ۹۹۹، روشن می‌شود).

\screenshotplaceholder{نمای کلی مدار \lr{main}: سه نمونه‌ی زنجیره‌شده‌ی \lr{bcd\_digit} به‌همراه نمایشگرهای هفت‌قطعه‌ای}

\subsubsection{مثال عددی}

\[
123 + 456 = 579
\]

با ورودی‌های $A100{=}1, A10{=}2, A1{=}3$ و $B100{=}4, B10{=}5, B1{=}6$، مدار باید خروجی $S100{=}5, S10{=}7, S1{=}9$ و $\text{Cout}{=}0$ تولید کند.

\screenshotplaceholder{اجرای مدار در \lr{Logisim} با ورودی‌های $123$ و $456$ و مشاهده‌ی خروجی $579$ روی نمایشگرها}

% ==================================================================
\section{صحت‌سنجی و نتایج تست}
% ==================================================================

هر سه مدار با بردار تست (\lr{.vec}) و از طریق خط فرمان \lr{Logisim} در حالت \lr{headless} (\lr{--test-vector}) تست شدند:

\begin{center}
\begin{tabular}{@{}lll@{}}
\toprule
\textbf{مدار} & \textbf{تعداد حالت تست‌شده} & \textbf{نتیجه} \\
\midrule
\lr{full\_adder} & ۸ (تمام حالت‌های ممکن $2^3$) & ۸ / ۸ \checkmark \\
\lr{bcd\_digit} & ۲۰۰ (تمام ترکیب‌های معتبر $A\times B\times C_{in}$) & ۲۰۰ / ۲۰۰ \checkmark \\
\lr{main} & ۶ (نمونه‌های \lr{smoke test}) & ۶ / ۶ \checkmark \\
\bottomrule
\end{tabular}
\end{center}

عدد ۲۰۰ برای \lr{bcd\_digit} از این‌جا می‌آید: $A \in \{0..9\}$ (۱۰ حالت) $\times\ B \in \{0..9\}$ (۱۰ حالت) $\times\ C_{in} \in \{0,1\}$ (۲ حالت) $= 200$ حالت  -  یعنی \textbf{تمام} ورودی‌های معتبر \lr{BCD} پوشش داده شده‌اند، نه فقط چند نمونه.

\screenshotplaceholder{خروجی خط فرمان تست بردار برای مدار \lr{bcd\_digit}: \lr{Passed: 200, Failed: 0}}

% ==================================================================
\section{بحث و تحلیل طراحی}
% ==================================================================

\subsection{چرا از قطعه‌ی آماده‌ی \lr{Comparator} استفاده نشد؟}
در طراحی‌های اولیه، برای تشخیص «آیا جمع خام از ۹ بزرگ‌تر است»، از یک قطعه‌ی آماده‌ی \lr{Comparator} در حالت \lr{unsigned} استفاده می‌شد. اما این قطعه از نظر سخت‌افزاری بیش از حد لازم است: با فرمول کمینه‌ی معادل \eqref{eq:corr}، همان تشخیص با تنها ۲ گیت (\lr{OR} و \lr{AND}) به‌جای یک مقایسه‌گر کامل ۴ بیتی قابل انجام است. این جایگزینی هم تعداد قطعات را کم می‌کند و هم مدار را ساده‌تر و قابل‌فهم‌تر می‌کند.

\subsection{چرا ماژول جمع‌کننده‌ی ۴ بیتی جداگانه حذف شد؟}
در نسخه‌ی قبلی طراحی، یک مدار واسط به نام \lr{adder4} (جمع‌کننده‌ی ۴ بیتی عمومی، ساخته‌شده از ۴ \lr{full\_adder}) وجود داشت که \lr{bcd\_digit} از \textbf{دو نمونه} آن استفاده می‌کرد. در نسخه‌ی نهایی، این لایه‌ی واسط حذف و ۸ نمونه‌ی \lr{full\_adder} \textbf{مستقیما} داخل \lr{bcd\_digit} قرار گرفتند. نتیجه: دقیقا ۳ ماژول در کل پروژه (\lr{full\_adder}, \lr{bcd\_digit}, \lr{main})، بدون هیچ لایه‌ی سلسله‌مراتبی اضافه.

\subsection{هزینه‌ی سخت‌افزاری}
هر رقم \lr{bcd\_digit} از ۸ تمام‌جمع‌کننده (۴۰ گیت پایه) به‌علاوه‌ی ۳ گیت تصحیح تشکیل شده است؛ در مجموع، مدار سه‌رقمی \lr{main} از ۲۴ تمام‌جمع‌کننده و ۹ گیت تصحیح ساخته شده  -  یک مبادله‌ی معمول میان سادگی طراحی سلسله‌مراتبی و تعداد قطعات.

% ==================================================================
\section{نتیجه‌گیری}
% ==================================================================

در این آزمایش، یک جمع‌کننده‌ی دهدهی سه‌رقمی به‌طور کامل و سلسله‌مراتبی طراحی و پیاده‌سازی شد. نکات کلیدی به‌دست‌آمده:

\begin{itemize}
  \item جمع دو رقم \lr{BCD} نیاز به یک مرحله‌ی تصحیح دارد چون جمع باینری خام می‌تواند تا ۱۹ برسد در حالی که هر رقم \lr{BCD} حداکثر ۹ است.
  \item رقم نقلی باینری معمولی برای تشخیص این وضعیت کافی نیست؛ باید بازه‌ی ۱۰ تا ۱۵ نیز جداگانه تشخیص داده شود.
  \item با یک فرمول بولی کمینه (فرمول \eqref{eq:corr})، این تشخیص تنها با ۳ گیت  -  بدون نیاز به مقایسه‌گر کامل  -  قابل پیاده‌سازی است.
  \item طراحی سلسله‌مراتبی (\lr{full\_adder} $\to$ \lr{bcd\_digit} $\to$ \lr{main}) امکان استفاده‌ی مجدد و تست مستقل هر سطح را فراهم می‌کند.
  \item مدار نهایی با ۲۰۸ حالت تست (۸ برای \lr{full\_adder} و ۲۰۰ برای \lr{bcd\_digit}) به‌طور کامل و بدون خطا تأیید شده است.
\end{itemize}

\end{document}
